\documentclass{article}

% packages
\usepackage{amsmath, amsthm, thmtools, amsfonts, amssymb, luacode, catchfile, tikzducks, hyperref, ifthen}
\ifcsname c@kobocompile\endcsname
	\usepackage[a5paper, total={1072pt, 1448pt}, margin=10pt, includeheadfoot]{geometry} % set page margins
\else
	\usepackage[a4paper, margin=50pt, includeheadfoot]{geometry}
\fi
\usepackage[shortlabels]{enumitem}
\usepackage[skip=3pt, indent=0pt]{parskip}

% language
\usepackage[bidi=basic, layout=tabular, provide=*]{babel}
\ifcsname c@english\endcsname
	\babelprovide[main, import]{english}
\else
	\babelprovide[main, import]{hebrew}
	\babelprovide{rl}
\fi
%\babelfont{rm}{Libertinus Serif}
\babelfont{rm}[Renderer=Harfbuzz]{Libertinus Serif}
\babelfont{sf}{Libertinus Sans}
\babelfont{tt}{Libertinus Mono}

% style
\AddToHook{cmd/section/before}{\clearpage}	% Add line break before section
\linespread{1.3}
\setcounter{secnumdepth}{0}		% Remove default number tags from sections, this won't do well with theorems
\AtBeginDocument{\setlength{\belowdisplayskip}{3pt}}
\AtBeginDocument{\setlength{\abovedisplayskip}{3pt}}
\graphicspath{ {../images/} }

% operators
\DeclareMathOperator\cis{cis}
\DeclareMathOperator\Sp{Sp}
\DeclareMathOperator\tr{tr}
\DeclareMathOperator\im{Im}
\DeclareMathOperator\re{Re}
\DeclareMathOperator\diag{diag}
\DeclareMathOperator*\lowlim{\underline{lim}}
\DeclareMathOperator*\uplim{\overline{lim}}
\DeclareMathOperator\rng{rng}
\DeclareMathOperator\Sym{Sym}
\DeclareMathOperator\Arg{Arg}
\DeclareMathOperator\Log{Log}
\DeclareMathOperator\dom{dom}
\DeclareMathOperator\supp{Supp}
\DeclareMathOperator\var{Var}
\DeclareMathOperator\cov{Cov}

% commands
%\renewcommand\qedsymbol{\textbf{מש''ל}}
%\renewcommand\qedsymbol{\fbox{\emoji{lizard}}}
\newcommand{\Aa}[0]{\mathcal{A}}
\newcommand{\Bb}[0]{\mathcal{B}}
\newcommand{\CC}[0]{\mathbb{C}}
\newcommand{\Cc}[0]{\mathcal{C}}
\newcommand{\EE}[0]{\mathbb{E}}
\newcommand{\FF}[0]{\mathbb{F}}
\newcommand{\Ff}[0]{\mathcal{F}}
\newcommand{\Ii}[0]{\mathcal{I}}
\newcommand{\Gg}[0]{\mathcal{G}}
\newcommand{\Ll}[0]{\mathcal{L}}
\newcommand{\Mm}[0]{\mathcal{M}}
\newcommand{\NN}[0]{\mathbb{N}}
\newcommand{\Nn}[0]{\mathcal{N}}
\newcommand{\PP}[0]{\mathbb{P}}
\newcommand{\Pp}[0]{\mathcal{P}}
\newcommand{\QQ}[0]{\mathbb{Q}}
\newcommand{\RR}[0]{\mathbb{R}}
\newcommand{\Rr}[0]{\mathcal{R}}
\newcommand{\Ss}[0]{\mathcal{S}}
\newcommand{\TT}[0]{\mathbb{T}}
\newcommand{\Uu}[0]{\mathcal{U}}
\newcommand{\Vv}[0]{\mathcal{V}}
\newcommand{\Ww}[0]{\mathcal{W}}
\newcommand{\ZZ}[0]{\mathbb{Z}}
\newcommand{\acts}[0]{\circlearrowright}
\newcommand{\explain}[2] {
	\begin{flalign*}
		 && \text{#2} && \text{#1}
	\end{flalign*}
}
\newcommand{\maketitleprint}[0]{ \begin{center}
	%\begin{tikzpicture}[scale=3]
	%	\duck[graduate=gray!20!black, tassel=red!70!black]
	%\end{tikzpicture}	
	\includegraphics[width=6cm]{cover}
\end{center}
}

% theorem commands
\newtheoremstyle{c_remark}
	{}	% Space above
	{}	% Space below
	{}% Body font
	{}	% Indent amount
	{\bfseries}	% Theorem head font
	{}	% Punctuation after theorem head
	{.5em}	% Space after theorem head
	{\thmname{#1}\thmnumber{ #2}\thmnote{ \normalfont{\text{(#3)}}}}	% head content
\newtheoremstyle{c_definition}
	{3pt}	% Space above
	{3pt}	% Space below
	{}% Body font
	{}	% Indent amount
	{\bfseries}	% Theorem head font
	{}	% Punctuation after theorem head
	{.5em}	% Space after theorem head
	{\thmname{#1}\thmnumber{ #2}\thmnote{ \normalfont{\text{(#3)}}}}	% head content
\newtheoremstyle{c_plain}
	{3pt}	% Space above
	{3pt}	% Space below
	{\itshape}% Body font
	{}	% Indent amount
	{\bfseries}	% Theorem head font
	{}	% Punctuation after theorem head
	{.5em}	% Space after theorem head
	{\thmname{#1}\thmnumber{ #2}\thmnote{ \text{(#3)}}}	% head content

\ifcsname c@english\endcsname
	\theoremstyle{plain}
	\newtheorem{theorem}{Theorem}[section]
	\newtheorem{lemma}[theorem]{Lemma}
	\newtheorem{proposition}[theorem]{Proposition}
	\newtheorem*{proposition*}{Proposition}
	%\newtheorem{corollary}[theorem]{אין חלופה עברית}

	\theoremstyle{definition}
	\newtheorem{definition}[theorem]{Definition}
	\newtheorem*{definition*}{Definition}
	\newtheorem{example}{Example}[section]
	\newtheorem{exercise}{Exercise}[section]

	\theoremstyle{remark}
	\newtheorem*{remark}{Remark}
	\newtheorem*{solution}{Solution}
	\newtheorem{conclusion}[theorem]{Conclusion}
	\newtheorem{notation}[theorem]{Notation}
\else
	\theoremstyle{c_plain}
	\newtheorem{theorem}{משפט}[section]
	\newtheorem{lemma}[theorem]{למה}
	\newtheorem{proposition}[theorem]{טענה}
	\newtheorem*{proposition*}{טענה}
	%\newtheorem{corollary}[theorem]{אין חלופה עברית}

	\theoremstyle{c_definition}
	\newtheorem{definition}[theorem]{הגדרה}
	\newtheorem*{definition*}{הגדרה}
	\newtheorem{example}{דוגמה}[section]
	\newtheorem{exercise}{תרגיל}[section]

	\theoremstyle{c_remark}
	\newtheorem*{remark}{הערה}
	\newtheorem*{solution}{פתרון}
	\newtheorem{conclusion}[theorem]{מסקנה}
	\newtheorem{notation}[theorem]{סימון}
\fi

% Questions related commands
\newcounter{question}
\setcounter{question}{1}
\newcounter{sub_question}
\setcounter{sub_question}{1}

\ifcsname c@english\endcsname
	\newcommand{\question}[1][0]{
		\ifthenelse{#1 = 0}{}{\setcounter{question}{#1}}
		\section{Question \arabic{question}}
		\addtocounter{question}{1}
		\setcounter{sub_question}{1}
	}

	\newcommand{\subquestion}[1][0]{
		\ifthenelse{#1 = 0}{}{\setcounter{sub_question}{#1}}
		\subsection{Part \alph{sub_question}}
		\addtocounter{sub_question}{1}
	}
\else
	\newcommand{\question}[1][0]{
		\ifthenelse{#1 = 0}{}{\setcounter{question}{#1}}
		\section{שאלה \arabic{question}}
		\addtocounter{question}{1}
		\setcounter{sub_question}{1}
	}

	\newcommand{\subquestion}[1][0]{
		\ifthenelse{#1 = 0}{}{\setcounter{sub_question}{#1}}
		\subsection{סעיף \localecounter{letters.gershayim}{sub_question}}
		\addtocounter{sub_question}{1}
	}
\fi

% import lua and start of document
\directlua{common = require ('../common')}

\GetEnv{AUTHOR}

% headers
\author{\AUTHOR}
\date\today

\title{פתרון ממ''ן 12 – חשבון אינפיניטסימלי 1 (20474)}

\begin{document}
\maketitle
\section{שאלה 1}
\subsection{סעיף א'}
\explain{נוכיח בלשון $\epsilon, N$ כי מתקיים}{$\lim_{n \to \infty} \sqrt{\frac{4n + 1}{n}} = 2$}
\begin{proof}
	\explain{נשים לב תחילה כי מתקיים}{
	$
		\frac{4n + 1}{n}
		= \frac{\frac{4n}{n} + \frac{1}{n}}{\frac{n}{n}}
		= \frac{4 + \frac{1}{n}}{1}
		= 4 + \frac{1}{n}
	$}
	לכל $n \ne 0$. לכן נוכל במקום להוכיח את קיום הגבול
	\[
		\lim_{n \to \infty} \sqrt{4 + \frac{1}{n}} = 2
	\]
	נגדיר $\epsilon > 0$ מספר ממשי. על־פי הגדרת הגבול צריך להתקיים:
	\[
		\left| \sqrt{4 + \frac{1}{n}} - 2 \right| < \epsilon
	\]
	תוכן השורש הוא תמיד לפחות 4, ולכן תוצאתו תמיד גדולה מ־2, בהתאם תוכן הערך המוחלט חיובי תמיד ומתקיים:
	\[
		\left| \sqrt{4 + \frac{1}{n}} - 2 \right|
		= \sqrt{4 + \frac{1}{n}} - 2 < \epsilon
	\]
	לכל $n > N$ כאשר $N \in \NN$. נשים לב כי
	\[
		{\left(2 + \sqrt{\frac{1}{n}}\right)}^2 = 4 + 2\sqrt{\frac{1}{n}} + \frac{1}{n} > 4 + \frac{1}{n} = {\left(\sqrt{4 + \frac{1}{n}}\right)}^2
		\rightarrow \sqrt{4 + \frac{1}{n}} < 2 + \sqrt{\frac{1}{n}}
	\]
	\explain{נגדיר כי צריך להתקיים} {$\sqrt{4 + \frac{1}{n}} - 2 < 2 + \sqrt{\frac{1}{n}} - 2 = \frac{1}{\sqrt{n}} < \epsilon$ }
	\explain{נגדיר}{ $N = \left\lceil \frac{1}{\epsilon^2} \right\rceil$ }
	במצב זה הגבול $(1)$ מתקיים לכל $\epsilon > 0$.
\end{proof}

\subsection{סעיף ב'}
\textbf{(i)}
יהיו $(a_n)$ סדרה ו־$L$ מספר ממשי. ננסח בלשון $\epsilon, N$ את הטענה
\[
	\lim_{n \to \infty} a_n \ne L
\]
תחילה נצרין את הטענה:
\[
	\lnot \forall \epsilon > 0 (
	\exists N \in \NN (
	\forall n > N (\left| a_n - L \right| < \epsilon)))
\]
נפשט את הפסוק:
\[
	\exists \epsilon > 0 (
	\forall N \in \NN (
	\exists n > N (\left| a_n - L \right| \ge \epsilon)))
\]
ננסח פסוק זה במילים: \\*
הטענה $\lim_{n \to \infty} a_n \ne L$ מתקיימת אם ורק אם
קיים $\epsilon > 0$ עבורו לכל $N$ טבעי קיים $n > N$
כך שמתקיים $\left| a_n - L \right| \ge \epsilon$. \\
\textbf{(ii)}
על־פי הגדרת התבדרות (2.13) סדרה נקראת מתבדרת כאשר אין מספר $L$
שעבורו מתקיים הגבול $\lim_{n \to \infty} a_n = L$.
ננסח את הטענה בעזרת הגדרת התכנסות בלשון $\epsilon, N$: \\*
סדרה $(a_n)$ היא מתבדרת אם ורק אם
לכל $L$ ממשי קיים $\epsilon > 0$ עבורו לכל $N$ טבעי קיים $n > N$
כך שמתקיים $\left| a_n - L \right| \ge \epsilon$.

\subsection{סעיף ג'}
נוכיח שהסדרה $(a_n)$ מתבדרת בלשון $\epsilon, N$, כאשר
\[
	a_n = \frac{{(-1)}^n n + 1}{n + 2}
\]
על־כן, נוכיח כי לכל $L$ ממשי שנבחר קיים $\epsilon > 0$ שעבורו לכל $N$ טבעי קיים $n > N$ שעבורו $\left| a_n - L \right| \ge \epsilon$.
\begin{proof}
	מתקיים על־פי חישוב ישיר
	\begin{align*}
		& n = 2 \rightarrow a_n = \frac{3}{4} \\
		& n = 3 \rightarrow a_n = -\frac{2}{5} \\
	\end{align*}
	נגדיר $\epsilon = \frac{1}{4}, N = 1$, ובשל כך הטענות תהינה תקפות לכל $n > 1$.
	אם $L \ge 1$ אז $|a_2 - L| > \epsilon$ והתנאי מתקיים.
	אילו $L \le 0$ אז $|a_3 - L| > \epsilon$ והתנאי מתקיים גם כן. \\*
	כאשר $0 < L < \frac{3}{4}$ נגדיר $\epsilon = a_2 - L$ ונראה כי התנאי מתקיים, שכן מספר זה מקיים את התנאי $\epsilon > 0$.
	כאשר $\frac{3}{4} \le L < 1$ אז נגדיר את $\epsilon = L - a_3$ ובהתאם התנאי עדיין יתקיים.
	בסך־הכול ראינו כי הסדרה מתבדרת לכל $L \in \RR$ ולכן מתבדרת לפי הגדרה ב'(ii). % chktex 36
\end{proof}

\section{שאלה 2}
נחשב את הגבולות הבאים, או נוכיח שאינם קיימים:

\subsection{סעיף א'}
\begin{align*}
	\lim_{n \to \infty} \sqrt{n^2 + {(-1)}^n} - n
	& = \lim_{n \to \infty} \frac{\left(\sqrt{n^2 + {(-1)}^n} + n\right)\left(\sqrt{n^2 + {(-1)}^n} - n\right)}{\sqrt{n^2 + {(-1)}^n} + n} \\
	& = \lim_{n \to \infty} \frac{n^2 + {(-1)}^n - n^2}{\sqrt{n^2 + {(-1)}^n} + n} \\
	& = \lim_{n \to \infty} \frac{{(-1)}^n}{\sqrt{n^2 + {(-1)}^n} + n} \\
\end{align*}
נוכל לראות כי מכנה הגבול חיובי וגדול מ־$2n - 1$ כמעט לכל $n$ ולכן המכנה שואף ל־$\infty$, כאשר המונה חסום ב־$1$,
לכן לפי משפט 2.22 והאריתמטיקה של הגבולות לקבוע כי
\[
	\lim_{n \to \infty} \sqrt{n^2 + {(-1)}^n} - n = 0
\]

\subsection{סעיף ב'}
\begin{align*}
	\lim_{n \to \infty} \frac{3n^3 - 2n^6 - 1}{n^4 - \pi n^5 + 5n}
	& = \lim_{n \to \infty} \frac{3n^3/n^6 - 2n^6/n^6 - 1/n^6}{n^4/n^6 - \pi n^5/n^6 + 5n/n^6} \\
	& = \lim_{n \to \infty} \frac{3/n^3 - 2 - 1/n^6}{1/n^2 - \pi /n + 5/n^5} \\
	& = \frac{-2}{0^+} & \text{על־פי אריתמטיקה של הגבולות} \\
	& = -\infty
\end{align*}
לכן הסדרה מתכנסת ל־$-\infty$ במובן הרחב.

\subsection{סעיף ג'}
\begin{align*}
	& \frac{\sqrt{3} n^2 - 1}{n^4} \le \frac{\lfloor \sqrt{3} n^2 \rfloor}{n^4} \le \frac{\sqrt{3} n^2}{n^4} \\
	& \frac{\sqrt{3}/n^2 - 1/n^4}{1} \le \frac{\lfloor \sqrt{3} n^2 \rfloor}{n^4} \le \frac{\sqrt{3}}{n^2} \\
\end{align*}
נראה כי גם כי כנביעה ממשפט 2.22
\[
	\lim_{n \to \infty} \left(\frac{\sqrt{3}}{n^2} - \frac{1}{n^4} \right)
	=
	\lim_{n \to \infty} \frac{\sqrt{3}}{n^2}
	= 0
\]
לכן לפי כלל הסנדוויץ' מתקיים גם
\[
	\lim_{n \to \infty} \frac{\lfloor \sqrt{3} n^2 \rfloor}{n^4} = 0
\]

\subsection{סעיף ד'}
נגדיר
\[
	a_n = \sqrt[n]{\frac{\displaystyle \prod_{i = 1}^{n} 2i - 1}{\displaystyle \prod_{i = 1}^{n} 2i }}
\]
מתקיים
\[
	a_n
	= \sqrt[n]{\displaystyle \prod_{i = 1}^{n} \frac{2i - 1}{2i}}
	\overset{(2)}{\le} \frac{\displaystyle \sum_{i = 1}^{n} \frac{2i - 1}{2i}}{n}
	= \displaystyle \sum_{i = 1}^{n} \frac{2i - 1}{2i \cdot n}
	\overset{(3)}{\le} \displaystyle n \cdot \frac{2n - 1}{2n \cdot n}
	= \frac{2n - 1}{2n} 
	= 1 - \frac{1}{2n} 
\]
\begin{enumerate}[(1)]
	\item נובע ממשפט אי־שוויון הממוצעים.
	\item לכל מספר בסכום נתאים מספר גדול או שווה ממנו, וניתן להוכיח כי $\frac{n - 1}{n} < \frac{n}{n + 1}$ לכל $n$ ולכן נבחר את המספר האחרון בסכום.
\end{enumerate}
ממשפט אי־שוויון הממוצעים נובע גם
\[
	a_n
	\ge \frac{n}{\sum_{i = 1}^n \frac{2i}{2i - 1}}
	\ge \frac{n}{n \cdot \frac{n2 - 1}{2n}}
	= \frac{2n}{2n - 1}
	= \frac{1}{1 - \frac{1}{2n}}
\]
עוד נראה כי
\[
	\lim_{n \to \infty} \frac{1}{1 - \frac{1}{2n}} = \lim_{n \to \infty} \left( 1 - \frac{1}{2n} \right) = 1
\]
לכן לפי כלל הסנדוויץ' $(a_n)$ מתכנסת ל־1.

\section{שאלה 3}
יהיו $(a_n)$ ו־$(b_n)$ סדרות כך שמתקיים $(1) \lim_{n \to \infty} a_n b_n = \infty$.

\subsection{סעיף א'}
נפריך את הטענה כי אם כמעט כל אברי $(a_n)$ ו־$(b_n)$ חיוביים אז אחת משתי הסדרות מתכנסת לאינסוף לפחות על־ידי דוגמה נגדית. \\*
נגדיר
\[
	a_n = \begin{cases}
		1 & \text{ זוגי} n \\
		n & \text{ אי־זוגי} n \\
	\end{cases},
	b_n = \begin{cases}
		n & \text{ זוגי} n \\
		1 & \text{ אי־זוגי} n \\
	\end{cases}
\]
ניתן לראות כי מכפלת הסדרות היא הסדרה $c_n = n$ אשר מקיימת את תנאי $(1)$, אך שתי הסדרות מתבדרות.

\subsection{סעיף ב'}
נוכיח כי אם כמעט כל אברי $(b_n)$ חיוביים אז גם כמעט כל אברי $(a_n)$ חיוביים.
\begin{proof}
	נגדיר כי $(b_n)$ חיובית לכל $n$ לפי משפט 2.44 בלא פגיעה בכלליות ההוכחה.
	בשל גבול $(1)$ ידוע כי לכל $M \in \RR$ קיים $N \in \NN$ כך שלכל $n > N$ מתקיים $a_n b_n > M$.
	ידוע כי $b_n$ חיובי וכי $M$ חיובי, אז גם $a_n$ חיובי אף הוא, שאם לא כן לא יתקיים השוויון $a_n b_n > M$.
	לכן $a_n$ חיובי לכל $n$.
\end{proof}

\subsection{סעיף ג'}
נפריך את הטענה כי $\lim_{n \to \infty} b_n \ne 0$ בעזרת דוגמה נגדית. \\*
נגדיר
\[
	a_n = n^3,
	b_n = \frac{1}{n}
\]
לכן
\[
	a_n b_n = n^2
\]
ובהתאם תנאי $(1)$ מתקיים. למרות זאת, לפי טענה 2.10 מתקיים $\lim_{n \to \infty} b_n = 0$ בסתירה לטענה.

\subsection{סעיף ד'}
נוכיח כי קיים $N \in \NN$ כך שלכל $n > N$ מתקיים $b_n \ne 0$. \\*
הוכחת טענה זו דומה ביותר להוכחת סעיף ב', אנו רואים כי $a_n b_n > M$, ולכן נוכל להניח גם $a_n b_n > 0$, לכן $a_n \ne 0 \land b_n \ne 0$.

\subsection{סעיף ה'}
נוכיח כי אם $\lim_{n \to \infty} b_n = 5$, אז $\lim_{n \to \infty} a_n = \infty$.
\begin{proof}
	נניח כי הגבול הראשון מתקיים והגבול השני לא בשלילה.
	אילו $a_n$ מתבדרת אז מכפלת הסדרות תתבדר גם היא, בניגוד ל־$(1)$, לכן נניח ש־$(a_n)$ מתכנסת.
	אילו הייתה הסדרה מתכנסת למספר סופי, אז לפי האריתמטיקה של הגבולות $a_n b_n$ הייתה מתכנסת אף היא למספר סופי בניגוד ל־$(1)$,
	לכן $(a_n)$ יכולה להתכנס ל־$\pm \infty$ בלבד.
	אנו יודעים כי $(b_n)$ חיובית לכמעט כל $n$, אילו הייתה $(a_n)$ מתכנסת ל־$-\infty$ אז היא הייתה שלילית לכמעט כל $n$.
	במקרה זה, מכפלת הסדרות הייתה שלילית לכמעט כל $n$, אך אנו יודעים כי היא חיובית על־פי גבול $(1)$ ולכן בסך־הכול $\lim_{n \to \infty} a_n = \infty$.
\end{proof}

\subsection{סעיף ו'}
נפריך את הטענה כי אם $b_n < a_n$ כמעט לכל $n$ אז $\lim_{n \to \infty} a_n = \infty$ בעזרת דוגמה נגדית: \\*
נגדיר
\[
	a_n = -n, b_n = -n - 1, a_n b_n = n (n + 1)
\]
במקרה זה גבול $(1)$ מתקיים ואף $b_n < a_n$ לכל $n$ אך $\lim_{n \to \infty} a_n = -\infty$ בסתירה לטענה.

\subsection{סעיף ז'}
נוכיח כי אם $0 < b_n < a_n$ כמעט לכל $n$ אז $\lim_{n \to \infty} a_n = \infty$.
\begin{proof}
	על־פי הגדרת שאיפה לאינסוף, לכל $M \in \RR$ קיים $N \in \NN$ כך שלכל $n > N$ מתקיים $a_n b_n > M$.
	לכן מתקיים גם $a_n a_n > a_n b_n > M$, אז $a_n > \sqrt{M}$. נגדיר $M' = \sqrt{M}$ ואנו רואים כי הגדרת השאיפה חלה על $(a_n)$.
	לכן
	\[
		\lim_{n \to \infty} a_n = \infty
	\]
\end{proof}

\end{document}
